\chapter{INTRODUCTION} \label{c1}
\justifying

\section{Background}

Communication is a fundamental human need, yet approximately 63 million deaf people in India face significant communication challenges in their daily interactions with the hearing world. Sign Language (ASL - American Sign Language) is their primary mode of communication, but most hearing individuals cannot understand it. This creates a persistent accessibility gap in educational institutions, healthcare facilities, government services, workplace environments, and public spaces \cite{kumar2020sign}.

The advent of computer vision and machine learning has opened new possibilities for bridging this communication gap. Real-time hand gesture recognition systems can automatically translate sign language into text or speech, enabling seamless communication between deaf and hearing individuals. This project implements such a system using modern computer vision techniques and machine learning algorithms.

\vspace{1em}
\section{Problem Statement}

Despite advances in technology, there remains no accessible, real-time solution for sign language interpretation that works on standard hardware. Current limitations include:

\begin{itemize}
    \item Manual interpretation is expensive and not always available
    \item Existing sign language recognition systems are either too complex or require specialized hardware
    \item Most systems have poor accuracy and cannot be used in real-time applications
    \item No affordable solution exists for everyday use by individuals
\end{itemize}

This project addresses these challenges by developing a real-time ASL gesture recognition system that enables seamless, real-time communication between deaf and hearing communities through accessible technology that anyone can use with just a standard webcam and laptop.

\vspace{1em}
\section{Objectives}

The main objectives of this project are:

\begin{enumerate}
    \item \textbf{Detect Hand Gestures}: Implement real-time hand detection using MediaPipe to extract 21 key landmarks per hand
    \item \textbf{Recognize ASL Alphabets}: Build a machine learning classifier to recognize all 26 ASL letters (A-Z)
    \item \textbf{Real-Time Processing}: Ensure the system processes video at 30 FPS with minimal latency
    \item \textbf{Text Output}: Buffer recognized letters into words and display on screen
    \item \textbf{Text-to-Speech}: Convert recognized text to speech for better accessibility
    \item \textbf{User Interface}: Create an intuitive web-based interface using Streamlit and Flask
\end{enumerate}

\vspace{1em}
\section{Scope}

This project encompasses the recognition of all 26 ASL alphabets (A--Z) through real-time video processing at 30 FPS. The system produces text output with optional text-to-speech functionality, accessible via a web-based user interface built using Streamlit. It is designed to operate on standard hardware without requiring a GPU, and supports single hand detection.


